
\textbf{Lemme 1.10.} \uline{Avec les notations et les hypothèses de} 1.1, \uline{soient} \(\underline{F}\) et \(\underline{G}\) \uline{deux faisceaux localement libres de type fini sur} \(X\), \uline{tels que} :
\[
  \cl_{K^{\bullet}(X)}(\underline{F})=\cl_{K^{\bullet}(X)}(\underline{G});
\]
\uline{alors il existe un entier} \(n_0\), \uline{tel que pour} \(n\geq n_0\), \uline{on ait} :
\[
\begin{array}{ll}
\text{i)}& f_*(\underline{F}(n)),\ f_*(\underline{G}(n))\text{ \uline{sont localement libres de type fini sur} }S;\\[1mm]
\text{ii)}& \cl_{K^{\bullet}(S)}\bigl(f_*(\underline{F}(n))\bigr)=\cl_{K^{\bullet}(S)}\bigl(f_*(\underline{G}(n))\bigr).
\end{array}
\]

Soit \(A\) le groupe abélien libre engendré par les classes (à isomorphisme près) de faisceaux localement libres de type fini sur \(X\). Dire que, dans \(K^{\bullet}(X)\), \(\underline{F}\) et \(\underline{G}\) ont même classe, c'est dire, d'après la construction du \(K^{\bullet}\), qu'il existe des suites exactes de \(\underline{O}_X\)-Modules localement libres de type fini :
\[
  0\longrightarrow \underline{L}_i\longrightarrow \underline{M}_i\longrightarrow \underline{N}_i\longrightarrow 0,
  \qquad i=1,\ldots,k,
\]
et des entiers \(n_i\), tels que dans \(A\) on ait l'égalité :
\[
  (\underline{F})-(\underline{G})=\sum_i n_i\bigl((\underline{L}_i)-(\underline{M}_i)+(\underline{N}_i)\bigr).
\]
Pour tout \(n\), la tensorisation par \(\underline{O}_X(n)\) définit un automorphisme de \(A\), d'où :
\[
  (\underline{F}(n))-(\underline{G}(n))=\sum_i \bigl((\underline{L}_i(n))-(\underline{M}_i(n))+(\underline{N}_i(n))\bigr),
\]
les suites \(0\to \underline{L}_i(n)\to \underline{M}_i(n)\to \underline{N}_i(n)\to0\) restant exactes.

Soient \(B\) le groupe abélien libre engendré par les classes (à isomorphisme près) de faisceaux quasi-cohérents sur \(S\), et \(A'\subset B\) le sous-groupe engendré par les classes de faisceaux localement libres de type fini sur \(S\). Par image directe, on définit un homomorphisme \(A\to B\), et on a donc encore :
